\scanpage{8}
\begin{center}( 3 )\end{center}

\textsc{Corollaire.} --- \textit{Si $A$ est une localité sur le corps $K$, telle que $K_a$ et $L_a$ soient séparables sur $K$, il y a équivalence entre :}

(i) \textit{$D_K(A)$ est libre;}

(ii) $A$ est régulier.

\textsc{Proposition 1.} --- \textit{Si $A$ domine un corps $K$, tel que $D_K(A)$ soit libre et de rang fini, si $K_b$ (resp. $L_a$) est séparable sur $K_a$ (resp. $L_b$) et si}
\[
\operatorname{rg}\mathfrak n-\operatorname{rg}\mathfrak m=\{K_b:K_a\}-\{L_b:L_a\},
\]
\textit{$D_A(B)$ est un $B$-module de dimension homologique finie $\leqq1$.}

\textsc{Corollaire.} --- \textit{Si $A$ est un anneau local régulier contenant un corps $K$, tel que le couple $(K,A)$ jouisse des propriétés énoncées dans le corollaire ci-dessus, $D_A(B)$ est de dimension homologique $\leqq1$ sous les hypothèses de la proposition 1.}

\textsc{Proposition 2.} --- \textit{Soit $A$ un anneau local noethérien contenant un corps $K$ tel que $\{L_a:K\}$ soit fini et que $L_a$ soit séparable sur $K$. Si $D_K(A)$ est libre, de rang $\dim A+\{L_a:K\}$, $A$ est local régulier.}

Notons, pour conclure, que dans de nombreuses questions (\textit{cf.} géométrie formelle) il vaut mieux séparer $D_A(B)$ en considérant
\[
D'_A(B)=\frac{D_A(B)}{\displaystyle\bigcap_{n\geqq0}\mathfrak n^n D_A(B)};
\]
la proposition 2 est encore vraie, lorsqu'on y substitue $D'_K(A)$ à $D_K(A)$.

(*) Séance du 14 janvier 1963.

\textsuperscript{(1)} \textsc{Berger} et \textsc{Kunz}, \textit{Math. Z.}, \textbf{77}, 1961, p. 314-338.

\textsuperscript{(2)} \textsc{Bourbaki}, \textit{Algèbre}, chap. V, § 8, ex. 1.

\textsuperscript{(3)} \textsc{Bourbaki}, \textit{Algèbre comm.}, chap. I, § 2, ex. 17.

\textsuperscript{(4)} \textsc{Godement}, \textit{Séminaire Cartan-Chevalley}, 1955-1956, exposés 17 et 18.

\textsuperscript{(5)} \textsc{Nakai}, \textit{J. Math. Soc. Japan}, 1961.

\textsuperscript{(6)} \textsc{Nakai} et \textsc{Suzuki}, \textit{J. Sc. Hiroshima Univ.}, série A, \textbf{24}, n° 3, décembre 1960.

\hfill (\textit{École Normale Supérieure, 45, rue d'Ulm, Paris}, 5\textsuperscript{e}.)

\mspage{7}
